\documentclass[11pt,reqno]{amsart}
\usepackage[T1]{fontenc}
\usepackage{amsmath,amssymb,amsthm}
\usepackage[hidelinks]{hyperref}
\hypersetup{pdftitle={The endpoint constant for an extremal iid sum inequality},pdfauthor={Peiwei Wu and Qian Peng}}
\newtheorem{theorem}{Theorem}[section]
\newtheorem{proposition}[theorem]{Proposition}
\newtheorem{lemma}[theorem]{Lemma}
\newtheorem{corollary}[theorem]{Corollary}
\newcommand{\Prob}{\mathbb P}
\newcommand{\EE}{\mathbb E}
\title[The endpoint constant for an extremal iid sum inequality]{The endpoint constant for an extremal iid sum inequality}
\author{Peiwei Wu}
\author{Qian Peng}
\thanks{Qian Peng is the corresponding author.}
\date{November 28, 2026}
\subjclass[2020]{60E15}
\keywords{Extremal probability, iid linear inequalities, optimal discrete distributions, multinacci numbers}
\begin{document}
\begin{abstract}
For three independent nonnegative random variables with a common law, let
$S_2(c)$ be the supremum of the probability that the first two sum to
less than $c$ times the third. We prove that $2/3-S_2(c)$ is asymptotic
to $4\log c/[9\log(1/(2-c))]$ as $c$ increases to two.
The upper estimate for $S_2$ counts losses across bands below a quantile
of an arbitrary law. Optimizing the masses of a finite ladder gives the
matching estimate. At the level of our upper bound for the $m$th
multinacci coefficient, a single band estimate covers a run of
$\Theta(\sqrt{m\log m})$ consecutive coefficients.
Splitting each ladder site into two points improves on the leading approximation by a positive
multiple of $\log(m)/m^2$. Consequently a finite block whose
right-hand score exceeds the value at the $m$th coefficient requires
at least $m+(9/64)\log(m)-O(1)$ sites.
\end{abstract}
\maketitle

\section{The extremal probability}
Let $X,Y,Z$ be independent with common Borel probability law $\mu$ on
$[0,\infty)$. For $1<c<2$ write
\[
 P_c(\mu)=\Prob(X+Y<cZ),\qquad S_2(c)=\sup_\mu P_c(\mu),
 \qquad L(c)=\frac{\log(1/(2-c))}{\log c}.
\]
At least one member of every triple, its smallest, loses when chosen as
target. Thus $S_2(c)\le 2/3$. Extremal probabilities for iid linear
inequalities were studied by Bellec and Fritz~\cite{bf}. Zaytsev and
Bellec~\cite[\S\S6.1, 7.1.3]{zb} certified eighteen discontinuities of
$S_2$, constructed families of finite ladders, and asked whether the
function has infinitely many discontinuities. The bounds of~\cite{wp}
placed the limiting lower and upper values of $L(c)(2/3-S_2(c))$ between
$1/3$ and $1/2$. Our first result identifies the constant.

\begin{theorem}\label{th:limit}
For the supremum over all Borel probability laws on $[0,\infty)$,
\[
 \lim_{c\uparrow2}L(c)\bigl(2/3-S_2(c)\bigr)=\frac49.
\]
\end{theorem}

For $m\ge2$ let $c_m$ be the solution in $(1,2)$ of
$(2-c_m)c_m^m=1$. Equivalently, $\sum_{j=1}^m c_m^{-j}=1$.
This sum decreases strictly in its argument and increases in $m$.
Consequently $c_m$ is unique and $c_m\uparrow2$.
These multinacci coefficients are the ordinary ladder thresholds
of~\cite[\S6.1.7]{zb}. They satisfy $L(c_m)=m$.

\begin{theorem}\label{th:corr}
For all sufficiently large integers $m$,
\[
 S_2(c_m)\ \ge\ \frac23-\frac{4}{9m}
            +\frac{\log(m-3)}{1000m^2}.
\]
\end{theorem}

Section~\ref{sec:band} bounds $S_2$ uniformly on a right neighborhood
of $c_m$ containing a growing number of consecutive multinacci
coefficients. Section~\ref{sec:split}
constructs finite laws for Theorem~\ref{th:corr}.
Section~\ref{sec:blocks} uses the gain from splitting the sites to
bound the number of sites in a block that improves on $S_2(c_m)$.

\section{A uniform upper bound from quantile bands}\label{sec:band}
For a sample of three values let $J$ count the target
positions for which the other two values sum to less than $c$ times
the target. The variable $2-J$ is nonnegative and
$\EE(2-J)=3(2/3-P_c(\mu))$. This identity lets us count losses on
selected disjoint events.

We first recall a bound from~\cite{wp}. If $1<e<2$, then
\begin{equation}\label{eq:rough}
 2/3-S_2(e)\ge\frac1{3\lceil L(e)\rceil}.
\end{equation}
Indeed, put $h=\lceil L(e)\rceil$, $u_0=0$, and
$u_{i+1}=(1+u_i)/e$. Then $u_h=(1-e^{-h})/(e-1)\ge1$.
For an atomless law, condition on the maximum $z$ and its index.
In their original coordinate order the other two draws are independent
with law truncated below $z$. Divide $[0,z]$ at the $h-1$ interior
cutoffs $u_i z$. If the two draws $x\le y$ are in the same interval,
then $ey\le e u_{i+1}z=(1+u_i)z\le x+z$, and the middle target loses.
The conditional probability of a shared interval is at least $1/h$,
by the sum of squared interval probabilities. This proves
\eqref{eq:rough} for atomless laws. Adding independent uniform noise on
$[0,\delta]$ to each draw makes any law atomless. Under the coupling
$X\mapsto X+\delta U$ with $U$ uniform on $[0,1]$, every strict win
persists for all sufficiently small $\delta$. Fatou's lemma transfers
the estimate back to the original law.

\begin{proposition}\label{prop:band}
Suppose $r>1$ and $d,K\ge1$ are integers such that
$r(1+r^{-d})\le c<2$ and $r^{-(K+d)}\le 2-c$. Then
\begin{equation}\label{eq:band}
 S_2(c)\le \frac23-\frac4{9(\sqrt K+4\sqrt d)^2}.
\end{equation}
\end{proposition}
\begin{proof}
Take $\mu$ atomless, write $H=2/3-P_c(\mu)>0$ and set $s=2-c$.
If $dH\ge1$ the inequality follows immediately.
Otherwise let $\varepsilon=(dH)^{1/3}$, and choose $z>0$ with
$\mu([0,z])=1-\varepsilon$. For $i\ge0$ put
\[
 t_i=1-r^{-i},\qquad T_i=\mu([t_i z,\infty)),\qquad a_i=T_i-T_{i+1}.
\]
Here $T_0=1$ and $T_i\ge\varepsilon$. Since
\[
 t_i+t_{i+d}-ct_{i+1}
   =s+r^{-i}(c/r-1-r^{-d})\ge0,
\]
a triple with two values in $[t_i z,t_{i+1}z)$ and the third at least
$t_{i+d}z$ loses at its middle target. These events are disjoint as
$i$ varies. For every $n$ their contribution to $\EE(2-J)$ gives
\begin{equation}\label{eq:disjoint}
 H\ge\sum_{i<n}a_i^2 T_{i+d}\ge\varepsilon\sum_{i<n}a_i^2.
\end{equation}
All unqualified band sums start at $i=0$.

The interval $[t_K z,z]$ has mass $q=T_K-\varepsilon$. If $q>0$,
condition all three draws to lie there and rescale each as
$U=(X/z-1+\eta)/\eta$, with $\eta=r^{-K}$. A win implies
\[
 U_1+U_2<cU_3-s(1-\eta)/\eta
       \le (2-s/\eta)U_3\le (2-r^{-d})U_3 .
\]
The first weak inequality uses $U_3\le1$, and the second uses
$r^{-(K+d)}\le s$. For $e=2-r^{-d}$ we have $e>r$ since
$r+r^{1-d}<2$ and $r>1$. Thus $L(e)=d\log r/\log e\le d$.
By \eqref{eq:rough} and nonnegativity of $2-J$ off this event,
\begin{equation}\label{eq:tail}
 q^3\le 3(d+1)H\le6dH .
\end{equation}
This also holds when $q=0$.

We estimate the difference produced by the shift of $d$ indices.
The square-root difference formula, $2uv\le u^2+v^2$, and
\eqref{eq:disjoint} give
\begin{align*}
 \sum_{i<K}a_i(\sqrt{T_i}-\sqrt{T_{i+d}})
 &\le \frac1{2\sqrt\varepsilon}
       \sum_{j=0}^{d-1}\sum_{i<K}a_i a_{i+j} \\
 &\le \frac d{2\sqrt\varepsilon}\sum_{i<K+d} a_i^2
 \le \frac{dH}{2\varepsilon^{3/2}}.
\end{align*}
By Cauchy--Schwarz, $\sum_{i<K}a_i\sqrt{T_{i+d}}\le\sqrt{KH}$.
Using upper sums for the integral of $\sqrt t$ yields
\[
 \tfrac23(1-T_K^{3/2})=\int_{T_K}^1\sqrt t\,dt
 \le\sum_{i<K}a_i\sqrt{T_i}\le\sqrt{KH}+\tfrac12\sqrt{dH}.
\]
By \eqref{eq:tail}, $T_K\le(1+6^{1/3})(dH)^{1/3}$.
Since $(2/3)(1+6^{1/3})^{3/2}+1/2<4$, we get
$2/3\le(\sqrt K+4\sqrt d)\sqrt H$. Smoothing as above and taking
the supremum over $\mu$ proves \eqref{eq:band}.
\end{proof}

\begin{corollary}\label{cor:explicit}
For $m\ge32$ set $D_m=1+\log(4m)/\log(3/2)$. Then
\[
 S_2(c_m)\le C_m:=\frac23-
       \frac4{9(\sqrt{m+2}+4\sqrt{D_m})^2}.
\]
\end{corollary}
\begin{proof}
Take $r=2-1/m$, $d=\lceil\log(4m)/\log(3/2)\rceil$ and $K=m+2-d$.
Since $(3/2)^m\ge4m$ for $m\ge32$, we have $1\le d\le m$,
$r^{-d}\le1/(4m)$, and $2-c_m=c_m^{-m}<(3/2)^{-m}\le1/(4m)$.
Here $c_m\ge c_2=(1+\sqrt5)/2>3/2$. Thus
$r(1+r^{-d})\le2-1/(2m)\le c_m$.
Bernoulli's inequality gives $(r/2)^m\ge1/2$. Since $r^2\ge2$,
$r^{K+d}=r^{m+2}\ge2^m\ge c_m^m$. Apply
Proposition~\ref{prop:band} using $K\le m+2$ and $d\le D_m$.
\end{proof}

\begin{corollary}\label{cor:interval}
For $m\ge32$ let $d_m=\lceil\log(4m)/\log(3/2)\rceil$ and set
\[
 B_m=\frac23-\frac4{9(\sqrt{m+2}+4\sqrt{d_m})^2}.
\]
Then $B_m<C_m$, and $S_2(x)\le B_m$ throughout
$[c_m,2-(2-c_m)/(4m)]$.
\end{corollary}
\begin{proof}
Keep $r=2-1/m$ from Corollary~\ref{cor:explicit} and take $K=m+2$.
We have $r(1+r^{-d_m})\le c_m$ and
$r^{-(K+d_m)}\le(2-c_m)/(4m)$. Proposition~\ref{prop:band}
applies on the stated interval, and $d_m<D_m$ gives $B_m<C_m$.
\end{proof}

\begin{proposition}\label{prop:range}
Let $m\ge32$ and let $j\ge0$ be an integer. If $4j^2\le m d_m$, then
\[
 S_2(x)\le B_m\quad(c_m\le x\le2-2^{-(m+j)}),\qquad
 S_2(c_{m+j}+)\le B_m .
\]
Conversely, if Proposition~\ref{prop:band}
applies at $c_{m+j}$ with a bound at most $C_m$, then
\[
 j<2+8\sqrt{(m+2)D_m}+16D_m .
\]
The largest offset $j$ attainable at this level has order
$\sqrt{m\log m}$, even if the bound is evaluated at a larger
coefficient and transferred to $c_{m+j}$ by monotonicity.
\end{proposition}
\begin{proof}
Write $d=d_m$ and take $r=2-1/m$, $e=\lceil3d/4\rceil$ and
$K=m+2+j$. Since $8\le d\le m$, we have $e\le7d/8\le(15/16)^2d$.
Also $r\ge7/4$ and $r^3\ge(3/2)^4$, whence $r^e\ge(3/2)^d\ge4m$.
It follows that $r(1+r^{-e})\le2-1/(2m)\le c_m$. Using $j\le m$
and the estimates $(r/2)^m\ge1/2$ and $r^{m+2}\ge2^m$ from
Corollary~\ref{cor:explicit}, we obtain $r^{K+e}\ge2^{m+j}$.
The choice of $j$ gives
\[
 \sqrt K-\sqrt{m+2}\le\frac{j}{2\sqrt{m+2}}\le\frac{\sqrt d}{4},
 \qquad 4\sqrt e\le\frac{15}{4}\sqrt d .
\]
The resulting bound in Proposition~\ref{prop:band} is at most $B_m$.
Since $2-c_{m+j}=c_{m+j}^{-(m+j)}>2^{-(m+j)}$, the interval
contains a right neighborhood of $c_{m+j}$.

For the converse, take any admissible parameters $r,K,e$ at
$u=c_{m+j}$. The width condition gives $r^{K+e}\ge u^{m+j}$,
and the other condition gives $r<u$. Hence $m+j<K+e$.
A bound at most $C_m$ requires
$\sqrt K+4\sqrt e\le\sqrt{m+2}+4\sqrt{D_m}$. Squaring and using
$K+e\le(\sqrt K+4\sqrt e)^2$ proves the displayed upper bound.
If the estimate is taken at $v\in[c_n,2)$ with $n=m+j$, then $r<v<2$
and $r^{K+e}\ge(2-v)^{-1}\ge c_n^n$. Corollary~\ref{cor:explicit}
gives $c_n^n\ge(2-1/n)^n\ge2^{n-1}$, so $n\le K+e$ by integrality.
This yields the same upper bound with a weak inequality. The order
claim follows since $d_m$ and $D_m$ are both asymptotic to
$\log m/\log(3/2)$.
\end{proof}

\section{Optimal masses on a finite ladder}
Give $N$ ordered sites probabilities $p_1,\ldots,p_N$.
For three independent site indices, call the ordered triple $(i,j,k)$
a rank win if $\min(i,j)<k$. Write $P_N$ for the maximum probability
of a rank win over the closed simplex of masses. Its objective is
$\sum_{k=1}^N p_k(1-(\sum_{j=k}^N p_j)^2)$.
For any increasing positive sites and $c<2$, a weak win
$y_i+y_j\le cy_k$ implies a rank win.

Conversely, suppose $1<e<2$ and $(2-e)e^N\le1$. Take
$y_k=\sum_{j=1}^k e^{-j}$ for $1\le k\le N$. Then $y_N\le1$,
$e y_1=1$, and $1+y_{k-1}=e y_k$ for $k\ge2$.
Every pair satisfies $y_i\le e y_k$, and every rank win satisfies
$y_i+y_j\le e y_k$. Subtracting the same positive number less than
$y_1$ from all sites makes every designated triple comparison strict,
because its two sides decrease by $2$ and $e$ times that number,
respectively. In particular,
\begin{equation}\label{eq:ranklower}
 S_2(c_m)\ge P_m.
\end{equation}

\begin{lemma}\label{lem:rec}
Put $x_N=2-3P_N$ and $t_N=\sqrt{1+x_N}$. Then $x_1=2$,
$x_N>0$, and
\begin{align}
 x_{N+1}&=\frac{2x_N}{t_N(1+t_N)},\label{eq:step}\\
 \frac1{x_{N+1}}-\frac1{x_N}
  &=\frac{t_N+2}{2(t_N+1)}
   =\frac34-\frac{x_N}{4(1+t_N)^2}.\label{eq:inverse}
\end{align}
In particular, $Nx_N\longrightarrow 4/3$.
\end{lemma}
\begin{proof}
Separate the smallest site of an $(N+1)$-site law, leaving mass $b$
above it. When the target lies above the smallest site and at least one
input lies at that site, the triple wins. Triples wholly above it
contribute at most $b^3P_N$, with equality at the optimal conditional law.
Thus $P_1=0$ and
$P_{N+1}=\max_{0\le b\le1}\{b-(1-P_N)b^3\}$.
Given $x_N>0$, differentiation yields the unique maximizer
$b=1/t_N$ and $P_{N+1}=2/(3t_N)$. This gives positivity and
\eqref{eq:step} by induction, and rearrangement gives \eqref{eq:inverse}.
The reciprocal increments lie between $1/2$ and $3/4$, so $x_N\to0$.
Their limit is $3/4$. Taking arithmetic means proves the assertion.
\end{proof}

\begin{proof}[Proof of Theorem~\ref{th:limit}]
At $c=c_m$, Corollary~\ref{cor:explicit} gives
$\liminf_m m(2/3-S_2(c_m))\ge4/9$ since $D_m=O(\log m)$.
Equation~\eqref{eq:ranklower} and Lemma~\ref{lem:rec} give the
matching upper limit. For $c_m\le c\le c_{m+1}$, comparison of
$\sum_{j=1}^m c^{-j}$ and $\sum_{j=1}^{m+1}c^{-j}$ with one gives
$m\le L(c)\le m+1$. Monotonicity of $S_2$ now squeezes the product
between $m(2/3-S_2(c_{m+1}))$ and
$(m+1)(2/3-S_2(c_m))$. Both tend to $4/9$.
\end{proof}

\section{Splitting each site into two points}\label{sec:split}
The rank model assigns zero probability of success to triples sampled
from a single site. We use two nearby points at each site to recover
part of this loss while preserving every rank win. The amount
recovered depends on the cubic sum of the optimal masses.

\begin{lemma}\label{lem:cubic}
There are maximizing masses $p_1,\ldots,p_N$ for $P_N$ with
\[
 \sum_{i=1}^N p_i^3\ge \tfrac18 x_N^2 U_N,\qquad
 U_N=\sum_{k=1}^N x_k\ge\log(N+1).
\]
\end{lemma}
\begin{proof}
At stage $j\ge2$ of the recursion in Lemma~\ref{lem:rec}, the
previous sites retain the fraction $b_j=1/t_{j-1}$ and the newly
inserted smallest site receives $a_j=1-b_j=x_j/2$. The latter formula
also holds for the initial site, with $a_1=1$. Numbering sites by the
stage of insertion, their masses are $a_k\prod_{j=k+1}^N b_j$.
By \eqref{eq:step}, the inequality $b_j^3\ge(x_j/x_{j-1})^2$
is equivalent to $(1+t_{j-1})^2\ge4t_{j-1}$. Multiplying these
inequalities bounds the cube of the $k$th mass below by
$x_N^2 x_k/8$. Finally \eqref{eq:inverse} implies
$1/x_k\le1/2+3(k-1)/4\le k$. Summing $x_k\ge1/k$ yields the claim.
\end{proof}

\begin{proposition}\label{prop:split}
With $N=m-4$, for all sufficiently large $m$,
\begin{equation}\label{eq:split}
 S_2(c_m)\ge P_N+\tfrac3{64}x_N^2 U_N .
\end{equation}
\end{proposition}
\begin{proof}
Put $c=c_m$, $s=2-c$, $e=2-8s$, and $a=1-2s$.
For large $m$ we have $0<a<1$, $1<e<ca$ and $c^4\ge8$. Hence
$(2-e)e^N\le8s c^{m-4}=8/c^4\le1$.
Use the $N$ sites realizing the weak rank comparisons at $e$ from
Section~3, with maximizing masses $p_i$. Conditional on site $y_i$,
give each of $ay_i$ and $y_i$ probability $1/2$, independently for
each draw. Every coarse rank win survives strictly, since its perturbed
input sum is at most $ey_k<cay_k$. When all three sites coincide,
a high target with at least one low input wins, as $1+a=2-2s<c$.
The conditional success probability of these events is $3/8$.
Their coarse indices coincide, so they are disjoint from the rank wins.
We have $S_2(c_m)\ge P_N+(3/8)\sum_i p_i^3$. Apply Lemma~\ref{lem:cubic}.
\end{proof}

\begin{proof}[Proof of Theorem~\ref{th:corr}]
By \eqref{eq:inverse} and $1+t_k\ge2$,
\[
 \tfrac12+\tfrac34(N-1)-U_N/16
 \le 1/x_N \le \tfrac34 N.
\]
For $N=m-4$, this implies
$0\le 3m/4-1/x_N\le 4+U_N/16$. Because $4/(3m)\le x_N$,
\[
 0\le x_N-\frac4{3m}
     =x_N\frac4{3m}\left(\frac{3m}4-\frac1{x_N}\right)
     \le x_N^2(4+U_N/16).
\]
In \eqref{eq:split} substitute $P_N=2/3-x_N/3$. The result is
\[
 S_2(c_m)-\frac23+\frac4{9m}
 \ge x_N^2\left(\frac5{192}U_N-\frac43\right).
\]
Now $U_N\ge\log(m-3)$ and $x_N\ge1/m$. For large $m$ the
parentheses exceed $\log(m-3)/1000$, as required.
\end{proof}

\section{A lower bound on the size of an improving block}\label{sec:blocks}
Let $\nu$ be a probability law on $\ell$ positive sites. For independent
$X,Y,Z$ with law $\nu$, define
\[
 B_c(\nu)=\Prob(X+Y\le cZ),\qquad Q_c(\nu)=\Prob(X\le cZ).
\]
For $r\ge0$ the weak block score is
\[
 W_c(\nu,r)=\frac{B_c(\nu)+2rQ_c(\nu)+r^2}{1+3r+3r^2}.
\]
To see the role of this expression, assign a lower-scale copy relative
mass $r$ against an upper block of mass one. At a coefficient just
above $c$, all the weak comparisons of a finite block become strict.
If the target lies in the upper block, the terms in the numerator
count guaranteed wins with zero, one, or two lower inputs. The
denominator $(1+r)^3-r^3$ removes triples wholly in the lower copy,
where the construction repeats. We call a block improving at $c$
when $W_c(\nu,r)>S_2(c)$.

Independently of the placement of the sites, $B_c(\nu)\le P_\ell$
by rank order, and $Q_c(\nu)\le1$. Thus
\begin{equation}\label{eq:ceiling}
 W_c(\nu,r)\le H_\ell,\qquad
 H_\ell=\sup_{v\ge0}\frac{P_\ell+2v+v^2}{1+3v+3v^2}.
\end{equation}
The complete weak rank pattern on $m$ sites realizes this upper
bound at $c_m$ after taking the supremum over $r$, since every pair
comparison holds there. For all sufficiently large $m$, the unrestricted
strict population exceeds this benchmark at $c_m$.

\begin{theorem}\label{th:labels}
For every fixed integer $k\ge0$, $S_2(c_m)> H_{m+k}$ for all
sufficiently large $m$. More quantitatively, for all sufficiently large
$m$, every improving block at $c_m$ has size
\[
 \ell>m-7+\tfrac9{64}\log(m-3).
\]
\end{theorem}
\begin{proof}
For large $j$, $x_j\le1/2$. The excess over $P_j$ of the fraction in
\eqref{eq:ceiling} has numerator $x_jv-(1-x_j)v^2$ and
denominator $1+3v+3v^2$. Bounding the positive part by the maximum
of the numerator yields $0\le H_j-P_j\le x_j^2/[4(1-x_j)]\le x_j^2$.
Equation~\eqref{eq:inverse} gives $x_i-x_{i+1}\le x_i^2$.
Since the $x_i$ decrease, for $j\ge N=m-4$ we obtain
\[
 H_j\le P_N+\bigl(1+(j-N)/3\bigr)x_N^2.
\]
For $j<N$ use $H_j\le H_N$, by monotonicity of $P_j$.
By Proposition~\ref{prop:split} the gain above $P_N$, divided by
$x_N^2$, is at least $(3/64)U_N\ge(3/64)\log(m-3)$. This proves
$S_2(c_m)>H_{m+k}$ for fixed $k$ and large $m$. If $\ell$ is the
size of an improving block, then $\ell>N$ for large $m$, and
\[
 (3/64)\log(m-3)<1+(\ell-N)/3.
\]
Rearrangement proves the stated bound.
\end{proof}

A jump at $c_m$ would require a strict separation of $S_2(c_m)$
from its right limit. Theorem~\ref{th:labels} excludes a growing
range of block sizes from providing such a separation. The existence
of infinitely many jumps of the unrestricted function remains open.

\section*{Formal proofs}
The accompanying Lean~4 development proves the endpoint limit,
the interval bounds and the estimates on the range of the band argument,
the positive correction, and both assertions on block sizes.
Its theorem index gives the declarations corresponding
to each result. In the formal rank recursion the index $n$ counts
one fewer than the number of sites. The source and build instructions
are available in the
\href{https://github.com/PEIWEIWU-AYS/one-scale-iid-probability-jumps-golden-ratio/releases/tag/endpoint-v1.1.0}{version 1.1.0 release}.
The permanent archive is
\href{https://doi.org/10.5281/zenodo.24614357}{doi:10.5281/zenodo.24614357}.

\begin{thebibliography}{9}
\raggedright
\bibitem{bf} P.~C. Bellec and T. Fritz,
On maximizing the probability of a linear inequality between iid random variables,
\emph{Amer. Math. Monthly} (2026),
\href{https://doi.org/10.1080/00029890.2026.2702871}{doi:10.1080/00029890.2026.2702871}.
\bibitem{wp} P. Wu and Q. Peng,
The jump set of the one-scale problem for iid linear inequalities,
preprint (2026),
\href{https://github.com/PEIWEIWU-AYS/one-scale-iid-probability-jumps-golden-ratio/releases/tag/v1.0.0}{version 1.0.0}.
\bibitem{zb} A. Zaytsev and P.~C. Bellec,
Flag algebras and the exact largest probability of a linear inequality between iid random variables,
\href{https://arxiv.org/abs/2611.29223}{arXiv:2611.29223} (2026).
\end{thebibliography}
\end{document}
